\documentclass[12pt,a4paper]{article}

% Encoding and language
\usepackage[T1]{fontenc}
\usepackage[utf8]{inputenc}
\usepackage[english]{babel}

% Page layout
\usepackage[margin=2.5cm]{geometry}

% Mathematics
\usepackage{amsmath,amssymb}

% Better code formatting
\usepackage{listings}

\title{Java and Spring Boot Interview Notes}
\author{İbrahim Utku Adanur}
\date{\today}

\begin{document}

\maketitle

\section{Java Interview Questions}

\subsection{What is the Difference Between \texttt{==} and \texttt{equals()}?}

In Java, \texttt{==} compares primitive values directly. For objects, it
checks whether two references point to the same object. The
\texttt{equals()} method is used for logical equality and can be overridden
by a class to compare object contents. The default implementation of
\texttt{equals()} in the \texttt{Object} class performs reference
comparison.

\subsubsection*{Key Point}

\begin{itemize}
    \item For primitives, \texttt{==} compares values.
    \item For objects, \texttt{==} compares references.
    \item \texttt{equals()} is used for logical or content equality.
\end{itemize}

\subsection{What Does the \texttt{final} Keyword Mean for Variables, Methods, and Classes?}

In Java, the \texttt{final} keyword prevents further modification depending
on where it is used. A final variable cannot be reassigned, a final method
cannot be overridden and a final class cannot be extended. For object
references, \texttt{final} prevents reassignment of the reference, but it
does not make the object itself immutable.

\subsection{What is the \texttt{@Qualifier} Annotation in Spring?}

The \texttt{@Qualifier} annotation is used when multiple beans of the same
type exist in the Spring ApplicationContext and Spring cannot determine
which one should be injected.

\subsubsection*{\texttt{@Qualifier} vs \texttt{@Primary}}

\texttt{@Primary} defines the default bean when multiple beans of the same
type exist.

\texttt{@Qualifier} explicitly selects a specific bean.

\begin{verbatim}
@Primary   -> preferred/default bean
@Qualifier -> explicitly selected bean
\end{verbatim}

\subsubsection*{Important Note}

Using \texttt{@Qualifier} binds that dependency to a specific implementation.

Therefore, it is suitable when the implementation is known when the
application is configured.

If the implementation must be selected dynamically at runtime, for example
based on a user's payment method, a Strategy or Factory-based design is
usually more appropriate.

\subsubsection*{Interview Answer}

The \texttt{@Qualifier} annotation is used when multiple beans of the same
type exist in the Spring context. It tells Spring exactly which bean should
be injected and resolves ambiguity during dependency injection.

It is mainly used for static dependency selection. If the implementation
must be chosen dynamically at runtime, a Strategy or Factory approach can
be used instead.

\subsection{Can a Static Method Be Overridden in Java?}

No. Static methods cannot be overridden in Java.

If a subclass defines a static method with the same signature as a static
method in its superclass, this is called \textit{method hiding}, not method
overriding.

\subsubsection*{Method Overriding vs Method Hiding}

\begin{verbatim}
Instance method:
-> Can be overridden
-> Runtime polymorphism
-> Actual object type determines the method

Static method:
-> Cannot be overridden
-> Can be hidden
-> Reference/class type determines the method
\end{verbatim}

\subsubsection*{Interview Answer}

Static methods cannot be overridden in Java because they belong to the
class rather than an instance. If a subclass defines a static method with
the same signature, it hides the superclass method. This is called method
hiding. Static method calls are resolved based on the reference or class
type rather than the actual object type.

\subsection{Which Memory Area is Cleaned by the Garbage Collector: Heap or Stack?}

The Java Garbage Collector primarily manages \textbf{heap memory}.

Objects created with \texttt{new} are generally allocated on the heap.

For example:

\begin{lstlisting}[language=Java]
Person person = new Person("Utku");
\end{lstlisting}

The local reference \texttt{person} may exist in the current stack frame,
while the actual \texttt{Person} object is stored on the heap.

If the object becomes unreachable, it becomes eligible for garbage
collection.

\begin{lstlisting}[language=Java]
Person person = new Person("Utku");

person = null;
\end{lstlisting}

After the reference is removed, if no other reachable reference points to
the object, the object may be reclaimed by the Garbage Collector.

Being eligible for garbage collection does not mean that the object is
immediately removed. The JVM decides when garbage collection should run.

\subsubsection*{What Happens to Stack Memory?}

Stack memory is managed automatically using stack frames.

Each method call creates a stack frame that contains information such as
local variables and references.

\begin{lstlisting}[language=Java]
void calculate() {
    int x = 10;
    int y = 20;
}
\end{lstlisting}

When the method finishes, its stack frame is automatically removed.
Therefore, the Garbage Collector does not normally clean stack memory.

\begin{verbatim}
Heap:
-> Stores objects
-> Managed mainly by the Garbage Collector
-> Unreachable objects may be reclaimed

Stack:
-> Stores method frames and local variables
-> Automatically managed
-> Frames are removed when methods return
\end{verbatim}

\subsubsection*{Interview Answer}

The Java Garbage Collector mainly manages heap memory. Objects are allocated
on the heap and when they are no longer reachable, they become eligible for
garbage collection. Stack memory is managed automatically through stack
frames, which are removed when methods return.

\subsection{Why is \texttt{String} Immutable in Java?}

A \texttt{String} object cannot be modified after it has been created.
Operations that appear to modify a String actually create a new String
object.

\subsubsection*{Interview Answer}

String is immutable in Java mainly for security, safe String Pool sharing,
thread safety, and reliable hashing. Since a String cannot change after
creation, it can safely be shared between threads, reused from the String
Pool, and used as a stable key in hash-based collections such as HashMap.

\subsection{What is the Difference Between \texttt{HashSet} and \texttt{TreeSet}?}

Both \texttt{HashSet} and \texttt{TreeSet} implement the \texttt{Set}
interface, meaning that they do not allow duplicate elements.

\begin{verbatim}
HashSet:
-> Unique elements
-> No guaranteed ordering
-> Hash-based
-> Uses hashCode() and equals()
-> Average O(1) operations

TreeSet:
-> Unique elements
-> Sorted elements
-> Tree-based
-> Uses Comparable or Comparator
-> O(log n) operations
\end{verbatim}

\subsubsection*{Interview Answer}

Both HashSet and TreeSet store unique elements. HashSet uses hashing and
does not guarantee ordering, while TreeSet keeps elements sorted using a
tree structure. HashSet provides average O(1) operations, whereas TreeSet
provides O(log n) operations. HashSet relies mainly on hashCode() and
equals(), while TreeSet relies on Comparable or Comparator.

\subsection{What is the Difference Between \texttt{HashMap} and \texttt{HashSet}?}

\texttt{HashMap} stores data as key-value pairs, while \texttt{HashSet}
stores only unique elements.

\subsubsection*{Interview Answer}

HashMap stores key-value pairs, while HashSet stores only unique elements.
HashMap does not allow duplicate keys, although duplicate values are
allowed. HashSet does not allow duplicate elements. Internally, HashSet is
backed by a HashMap, where the set elements are stored as map keys with a
dummy value.

\subsection{What is a Bean?}

A bean is simply an object whose lifecycle is managed by Spring.

Spring may create beans automatically by discovering classes annotated
with annotations such as:

\begin{verbatim}
@Component
@Service
@Repository
@Controller
@RestController
\end{verbatim}

Beans can also be explicitly declared using the \texttt{@Bean}
annotation.

\subsection{What is the ApplicationContext?}

The \texttt{ApplicationContext} is the central Spring IoC container.

It is responsible for creating, storing, configuring, and connecting
Spring-managed objects, which are called \textbf{beans}.

\texttt{ApplicationContext} can be thought of as the
container in which Spring-managed objects live.

\subsection{What Does \texttt{@SpringBootApplication} Do Internally?}

\texttt{@SpringBootApplication} is a convenience annotation that combines:

\begin{verbatim}
@SpringBootConfiguration
@EnableAutoConfiguration
@ComponentScan
\end{verbatim}

\texttt{@SpringBootConfiguration} is based on
\texttt{@Configuration} and allows Spring configuration and bean
definitions.

For example:

\begin{lstlisting}[language=Java]
@Configuration
public class AppConfig {

    @Bean
    public PaymentClient paymentClient() {
        return new PaymentClient();
    }
}
\end{lstlisting}

The object returned by a \texttt{@Bean} method is registered in the
Spring \texttt{ApplicationContext}.

\texttt{@ComponentScan} searches the application's package and its
subpackages for classes such as:

\begin{verbatim}
@Component
@Service
@Repository
@Controller
@RestController
\end{verbatim}

and registers them as Spring beans.

\texttt{@EnableAutoConfiguration} automatically configures the application
based on the libraries available on the classpath, application properties
and existing beans.

For example, if \texttt{spring-boot-starter-web} is present, Spring Boot
automatically configures much of the Spring MVC infrastructure.

Finally:

\begin{lstlisting}[language=Java]
SpringApplication.run(DemoApplication.class, args);
\end{lstlisting}

starts the application and initializes the Spring
\texttt{ApplicationContext}.

\subsubsection*{Interview Answer}

\texttt{@SpringBootApplication} combines configuration, component scanning
and auto-configuration. It allows Spring to discover and create beans,
configure the application automatically and manage those beans inside the
\texttt{ApplicationContext}. The \texttt{ApplicationContext} is the IoC
container responsible for managing beans and injecting their dependencies.

\subsection{What is \texttt{@Autowired} in Spring?}

\texttt{@Autowired} is used for dependency injection. It tells Spring to
find a matching bean from the \texttt{ApplicationContext} and inject it.

\begin{lstlisting}[language=Java]
@Service
public class OrderService {

    private final PaymentService paymentService;

    public OrderService(PaymentService paymentService) {
        this.paymentService = paymentService;
    }
}
\end{lstlisting}

Since Spring 4.3, if a class has only one constructor,
\texttt{@Autowired} is not required because Spring automatically uses that
constructor for dependency injection.

If a class has multiple constructors, \texttt{@Autowired} can be used to
indicate which constructor Spring should use.

\texttt{final} is commonly used with constructor injection because the
dependency is assigned once during object creation and cannot later be
reassigned to another object.

\texttt{@Autowired} can also be used for setter and field injection, but
constructor injection is generally preferred.

\subsubsection*{Interview Answer}

\texttt{@Autowired} tells Spring to inject a matching bean from the
\texttt{ApplicationContext}. With a single constructor, it is usually not
needed. With multiple constructors, it can specify which constructor should
be used for injection. Constructor-injected dependencies are often declared
\texttt{final} so their references cannot be reassigned.

\subsection{What is a Singleton?}

A singleton is a design pattern that ensures only one instance of a class
exists and provides a shared access point to that instance.

In Spring, the default bean scope is \texttt{singleton}. This means Spring
creates one bean instance per \texttt{ApplicationContext} and reuses that
same instance whenever the bean is injected.

For example:

\begin{lstlisting}[language=Java]
@Service
public class OrderService {
}
\end{lstlisting}

If \texttt{OrderService} is injected into multiple classes, Spring normally
provides the same bean instance.

\begin{verbatim}
ApplicationContext
        |
        +-- one OrderService instance
              |
              +-- used by Controller A
              +-- used by Service B
\end{verbatim}

A Spring singleton is not automatically thread-safe. If the bean contains
mutable shared state, multiple threads may access and modify the same data.

\subsubsection*{Interview Answer}

A singleton means that only one instance of an object is used within a
given context. In Spring, singleton is the default bean scope, so one bean
instance is created per ApplicationContext and shared wherever that bean is
injected. However, singleton does not automatically mean thread-safe.

\subsection{Are Singleton Beans Thread-Safe in Spring?}

No. A singleton bean is not automatically thread-safe.

By default, Spring creates one bean instance per
\texttt{ApplicationContext}. Multiple threads may use the same bean
instance at the same time.

A stateless singleton bean is generally safe:

\begin{lstlisting}[language=Java]
@Service
public class OrderService {

    public int calculateTotal(int price, int quantity) {
        return price * quantity;
    }
}
\end{lstlisting}

However, mutable shared state may cause race conditions:

\begin{lstlisting}[language=Java]
@Service
public class CounterService {

    private int counter = 0;

    public int increment() {
        return ++counter;
    }
}
\end{lstlisting}

Multiple threads can modify \texttt{counter} concurrently.

\subsubsection*{Interview Answer}

Spring singleton beans are not automatically thread-safe. A singleton means
that one bean instance is shared within the ApplicationContext. Stateless
beans are generally safe, but mutable shared state can cause race
conditions and must be handled carefully.

\subsection{What is the Difference Between Abstraction and Encapsulation?}

\subsubsection*{Interview Answer}

Abstraction hides implementation details and focuses on what an object does.
Encapsulation protects an object's internal state by controlling access to
its data. Abstraction is about hiding complexity, while encapsulation is
about restricting direct access to data.

\subsection{What is the Single Responsibility Principle?}

\subsubsection*{Interview Answer}

The Single Responsibility Principle states that a class should have one
main responsibility and only one reason to change. Unrelated responsibilities
should be separated into different classes.

\subsection{What is the Open/Closed Principle?}

\subsubsection*{Interview Answer}

The Open/Closed Principle states that software components should be open
for extension but closed for modification. New behavior should preferably
be added through new implementations instead of repeatedly changing
existing code.

\subsection{What is the Liskov Substitution Principle?}

The Liskov Substitution Principle states that a subclass should be usable
wherever its parent class is expected without breaking the expected behavior.

\begin{lstlisting}[language=Java]
class Vehicle {
    public void startEngine() {
        System.out.println("Engine started");
    }
}

class Bicycle extends Vehicle {

    @Override
    public void startEngine() {
        throw new UnsupportedOperationException();
    }
}
\end{lstlisting}

A \texttt{Bicycle} cannot safely replace a \texttt{Vehicle} because it
cannot provide the expected \texttt{startEngine()} behavior.

\subsubsection*{Interview Answer}

A subclass should be replaceable with its parent type without breaking the
expected behavior of the program.

\subsection{What is the Interface Segregation Principle?}

The Interface Segregation Principle states that a class should not be forced
to implement methods it does not need.

For example:

\begin{lstlisting}[language=Java]
interface Worker {
    void work();
    void eat();
}

class Robot implements Worker {

    public void work() {
        System.out.println("Working");
    }

    public void eat() {
        // Robot does not need this method
    }
}
\end{lstlisting}

A better design uses smaller and more focused interfaces:

\begin{lstlisting}[language=Java]
interface Workable {
    void work();
}

interface Eatable {
    void eat();
}

class Robot implements Workable {
    public void work() {
    }
}

class Human implements Workable, Eatable {
    public void work() {
    }

    public void eat() {
    }
}
\end{lstlisting}

\subsubsection*{Interview Answer}

The Interface Segregation Principle states that classes should not be forced
to depend on methods they do not use. Small and focused interfaces are
preferred over large general-purpose interfaces.

\subsection{What is the Dependency Inversion Principle?}

The Dependency Inversion Principle states that high-level classes should
depend on abstractions rather than concrete implementations.

For example, this creates tight coupling:

\begin{lstlisting}[language=Java]
class OrderService {

    private CreditCardPaymentService paymentService =
            new CreditCardPaymentService();

    public void checkout() {
        paymentService.pay();
    }
}
\end{lstlisting}

A better design depends on an abstraction:

\begin{lstlisting}[language=Java]
interface PaymentService {
    void pay();
}

class CreditCardPaymentService
        implements PaymentService {

    public void pay() {
        System.out.println("Paid by credit card");
    }
}

class OrderService {

    private final PaymentService paymentService;

    public OrderService(PaymentService paymentService) {
        this.paymentService = paymentService;
    }

    public void checkout() {
        paymentService.pay();
    }
}
\end{lstlisting}

Now \texttt{OrderService} depends on the
\texttt{PaymentService} abstraction instead of a specific payment
implementation.

\subsubsection*{Interview Answer}

The Dependency Inversion Principle means that high-level classes should
depend on abstractions rather than concrete implementations. For example,
\texttt{OrderService} should depend on \texttt{PaymentService}, not directly
on \texttt{CreditCardPaymentService}.

\subsection{How Can Objects Be Created Without Using the \texttt{new} Keyword?}

Apart from directly using the \texttt{new} keyword, Java objects can be
created or obtained in several ways.

\subsubsection*{Reflection}

Reflection can invoke a constructor dynamically at runtime.

\begin{lstlisting}[language=Java]
User user = User.class.getDeclaredConstructor().newInstance();
\end{lstlisting}

The older \texttt{Class.newInstance()} method is deprecated. Modern Java
prefers \texttt{getDeclaredConstructor().newInstance()}.

\subsubsection*{Cloning}

An object can be created as a copy of another object using
\texttt{clone()}.

\begin{lstlisting}[language=Java]
User user2 = (User) user1.clone();
\end{lstlisting}

The class normally needs to implement \texttt{Cloneable}.

\subsubsection*{Deserialization}

A previously serialized object can be reconstructed from a stream.

\begin{lstlisting}[language=Java]
ObjectInputStream in =
    new ObjectInputStream(
        new FileInputStream("user.ser")
    );

User user = (User) in.readObject();
\end{lstlisting}

\subsubsection*{Factory Methods}

Objects can also be obtained through factory methods:

\begin{lstlisting}[language=Java]
Integer number = Integer.valueOf(10);
\end{lstlisting}

A factory method may internally create a new object or return an existing
one.

\subsubsection*{Interview Answer}

Apart from using the \texttt{new} keyword directly, objects can be created
using reflection, cloning, and deserialization. Objects can also be obtained
through factory methods or frameworks such as Spring.

\subsection{What are the Different ClassLoaders in the JVM?}

A \texttt{ClassLoader} is responsible for loading Java classes into the JVM.

The main built-in class loaders are:

\begin{verbatim}
Bootstrap ClassLoader
Platform ClassLoader
Application ClassLoader
\end{verbatim}

\subsubsection*{Bootstrap ClassLoader}

The Bootstrap ClassLoader loads core Java classes such as
\texttt{java.lang.String} and \texttt{java.lang.Object}.

It is part of the JVM itself.

\subsubsection*{Platform ClassLoader}

The Platform ClassLoader loads Java platform-level classes.

In older Java versions, a similar role was performed by the
\texttt{Extension ClassLoader}.

\subsubsection*{Application ClassLoader}

The Application ClassLoader loads application classes and dependencies
available on the classpath.

For example, classes such as:

\begin{verbatim}
OrderService
JobApplicationController
PaymentService
\end{verbatim}

are normally loaded by the Application ClassLoader.

Class loaders usually follow the \textbf{parent delegation model}, where a
class loader first asks its parent to load the requested class.

\subsubsection*{Interview Answer}

The JVM mainly uses Bootstrap, Platform, and Application ClassLoaders.
Bootstrap loads core Java classes, Platform loads platform-level classes,
and Application loads application classes and classpath dependencies.
They normally follow the parent delegation model.

\subsection{What is the Difference Between Heap and Stack Memory in Java?}

\subsubsection*{Interview Answer}

Stack memory stores method frames and local variables, while heap memory
stores objects. Each thread has its own stack, whereas the heap is shared.
Stack frames are removed automatically when methods return, while heap
objects are managed by the Garbage Collector.

\subsection{What is the Difference Between an Exception and an Error in Java?}

Both \texttt{Exception} and \texttt{Error} inherit from
\texttt{Throwable}.

An \texttt{Exception} usually represents a problem that an application can
anticipate and potentially handle.

Examples include:

\begin{verbatim}
IOException
SQLException
NullPointerException
IllegalArgumentException
\end{verbatim}

An \texttt{Error} represents a more serious JVM or system-level problem.

Examples include:

\begin{verbatim}
OutOfMemoryError
StackOverflowError
\end{verbatim}

Although an \texttt{Error} can technically be caught, applications normally
should not try to recover from such serious failures.

\subsubsection*{Interview Answer}

Exceptions usually represent application-level problems that can potentially
be handled, while Errors represent serious JVM or system-level failures that
applications normally should not try to recover from.

\subsection{What is the Purpose of the \texttt{finally} Block in Java?}

The \texttt{finally} block is used for code that should execute whether an
exception occurs or not.

It is commonly used for cleanup operations such as releasing resources.

The \texttt{finally} block normally executes regardless of whether the
\texttt{try} block succeeds or throws an exception.

In modern Java, \texttt{try-with-resources} is often preferred for closing
resources automatically.

\subsubsection*{Interview Answer}

The \texttt{finally} block is used for code that should run whether an
exception occurs or not, typically for cleanup and resource release.

\subsection{Which Collection Maintains Insertion Order and Avoids Duplicates?}

\texttt{LinkedHashSet} maintains insertion order while also preventing
duplicate elements.

\begin{lstlisting}[language=Java]
Set<String> names = new LinkedHashSet<>();

names.add("Utku");
names.add("Ali");
names.add("Utku");
names.add("Mehmet");

System.out.println(names);
\end{lstlisting}

The output is:

\begin{verbatim}
[Utku, Ali, Mehmet]
\end{verbatim}

The second \texttt{"Utku"} is ignored because sets do not allow duplicates,
while the insertion order is preserved.

\subsubsection*{Interview Answer}

\texttt{LinkedHashSet} is used when unique elements are required while
preserving insertion order.

\subsection{How Does \texttt{ConcurrentHashMap} Handle Concurrency?}

\texttt{ConcurrentHashMap} provides thread-safe access while allowing
multiple threads to operate on the map concurrently.

Unlike a fully synchronized map, it does not lock the entire map for every
operation.

In modern Java:

\begin{itemize}
    \item Read operations such as \texttt{get()} are mostly non-blocking.
    \item Write operations use fine-grained synchronization and CAS.
    \item Different threads can update different buckets concurrently.
\end{itemize}

For example:

\begin{lstlisting}[language=Java]
ConcurrentHashMap<String, Integer> map =
        new ConcurrentHashMap<>();

map.put("A", 1);
map.put("B", 2);
\end{lstlisting}

Two threads may work on different parts of the map at the same time without
locking the whole structure.

\texttt{ConcurrentHashMap} does not allow \texttt{null} keys or values.

\subsubsection*{Interview Answer}

\texttt{ConcurrentHashMap} provides thread-safe access without locking the
entire map. Reads are mostly non-blocking, while updates use fine-grained
synchronization and CAS on individual buckets. This reduces contention and
allows multiple threads to operate concurrently.

\subsection{What is the Difference Between Shallow Copy and Deep Copy?}

A \textbf{shallow copy} creates a new outer object, but nested object
references are shared.

\begin{lstlisting}[language=Java]
Person p1 =
    new Person("Utku", new Address("Ankara"));

Person p2 =
    new Person(p1.name, p1.address);
\end{lstlisting}

Both objects refer to the same \texttt{Address} object.

Therefore:

\begin{lstlisting}[language=Java]
p2.address.city = "Istanbul";
\end{lstlisting}

also affects the address seen through \texttt{p1}.

A \textbf{deep copy} also creates separate copies of nested mutable objects.

\begin{lstlisting}[language=Java]
Person p2 = new Person(
    p1.name,
    new Address(p1.address.city)
);
\end{lstlisting}

Now the two \texttt{Person} objects have independent
\texttt{Address} objects.

\subsubsection*{Interview Answer}

A shallow copy creates a new outer object but shares references to nested
objects. A deep copy creates independent copies of both the outer object
and its nested mutable objects.

\subsection{What is a Functional Interface in Java?}

A functional interface is an interface that contains exactly one abstract
method.

Its main purpose is to represent a behavior that can be implemented using a
lambda expression.

For example:

\begin{lstlisting}[language=Java]
@FunctionalInterface
interface Calculator {
    int calculate(int a, int b);
}
\end{lstlisting}

Without lambda expressions, we could create a separate implementation:

\begin{lstlisting}[language=Java]
class AddCalculator implements Calculator {

    @Override
    public int calculate(int a, int b) {
        return a + b;
    }
}

Calculator calculator = new AddCalculator();
\end{lstlisting}

We could also use an anonymous class:

\begin{lstlisting}[language=Java]
Calculator calculator = new Calculator() {

    @Override
    public int calculate(int a, int b) {
        return a + b;
    }
};
\end{lstlisting}

Here, the interface itself is not instantiated. Java creates an anonymous
class that implements the interface and creates an instance of that class.

With a lambda expression, the same behavior can be written much more
concisely:

\begin{lstlisting}[language=Java]
Calculator calculator = (a, b) -> a + b;
\end{lstlisting}

The lambda expression provides the implementation of the single abstract
method in the functional interface.

Functional interfaces also allow behavior to be passed as a method
parameter:

\begin{lstlisting}[language=Java]
public void process(
        int a,
        int b,
        Calculator calculator
) {
    System.out.println(
        calculator.calculate(a, b)
    );
}

process(10, 5, (a, b) -> a + b);
process(10, 5, (a, b) -> a - b);
process(10, 5, (a, b) -> a * b);
\end{lstlisting}

The method remains the same, but the behavior passed to it can change.

Common functional interfaces in Java include:

\begin{verbatim}
Predicate<T>     -> input -> boolean
Function<T, R>  -> input -> output
Consumer<T>     -> input -> no return value
Supplier<T>     -> no input -> output
\end{verbatim}

\subsubsection*{Interview Answer}

A functional interface is an interface with exactly one abstract method.
It allows a behavior to be represented with a lambda expression and passed
around as a value, for example as a method parameter. Without lambdas, the
same behavior would normally require a separate implementation class or an
anonymous class.

\subsection{How Can You Create Your Own Functional Interface in Java?}

A custom functional interface is created by defining an interface with
exactly one abstract method.

\begin{lstlisting}[language=Java]
@FunctionalInterface
public interface Calculator {

    int calculate(int a, int b);
}
\end{lstlisting}

It can then be implemented using a lambda expression:

\begin{lstlisting}[language=Java]
Calculator add = (a, b) -> a + b;

int result = add.calculate(5, 3);
\end{lstlisting}

The \texttt{@FunctionalInterface} annotation is optional, but recommended.
It allows the compiler to verify that the interface contains only one
abstract method.

A functional interface may still contain \texttt{default} and
\texttt{static} methods.

\subsubsection*{Interview Answer}

To create a custom functional interface, define an interface with exactly
one abstract method. The \texttt{@FunctionalInterface} annotation is
optional but recommended because it lets the compiler enforce this rule.

\subsection{What is the Purpose of \texttt{Optional} in Java?}

\texttt{Optional<T>} is a container that represents a value that may or may
not be present.

It is commonly used as a return type when a method may not find a result.

\begin{lstlisting}[language=Java]
Optional<User> optionalUser =
        repository.findById(id);
\end{lstlisting}

Instead of manually checking for \texttt{null}, methods such as
\texttt{orElse()}, \texttt{orElseGet()}, and
\texttt{orElseThrow()} can be used.

For example:

\begin{lstlisting}[language=Java]
public User getUserById(Long id) {
    return repository.findById(id)
            .orElseThrow(() ->
                new ResourceNotFoundException(
                    "User not found"
                )
            );
}
\end{lstlisting}

If the user exists, \texttt{orElseThrow()} returns the \texttt{User}.
If the value is absent, it throws the specified exception.

\subsubsection*{Interview Answer}

\texttt{Optional} represents a value that may or may not be present. It is
commonly used for return values to express absence explicitly and reduce
manual null checks.

\subsection{What is the Difference Between Intermediate and Terminal Operations in Java Streams?}

Intermediate operations transform a stream and return another stream.

Examples include:

\begin{verbatim}
filter()
map()
sorted()
distinct()
limit()
\end{verbatim}

Terminal operations complete the stream pipeline and produce a result or
side effect.

Examples include:

\begin{verbatim}
toList()
forEach()
count()
reduce()
collect()
findFirst()
\end{verbatim}

For example:

\begin{lstlisting}[language=Java]
List<String> names = users.stream()
        .filter(user -> user.getAge() >= 18)
        .map(User::getName)
        .toList();
\end{lstlisting}

Here, \texttt{filter()} and \texttt{map()} are intermediate operations,
while \texttt{toList()} is a terminal operation.

Intermediate operations are usually lazy and are executed when a terminal
operation is called.

\subsubsection*{Interview Answer}

Intermediate operations return another stream and build the processing
pipeline. Terminal operations finish the pipeline and produce the final
result. Intermediate operations are generally lazy.

\subsection{How Can You Access Properties from \texttt{application.properties} in Spring?}

A property can be injected using the \texttt{@Value} annotation.

For example:

\begin{verbatim}
app.name=Job Tracker
\end{verbatim}

\begin{lstlisting}[language=Java]
@RestController
public class AppController {

    @Value("${app.name}")
    private String appName;
}
\end{lstlisting}

Spring reads \texttt{app.name} from the configuration and injects its value
into \texttt{appName}.

For multiple related properties, \texttt{@ConfigurationProperties} is
usually preferred because it groups configuration values into a dedicated
class.

\subsubsection*{Interview Answer}

Properties can be accessed using \texttt{@Value}. For multiple related
properties, \texttt{@ConfigurationProperties} provides a cleaner and more
maintainable approach.

\subsection{What is the Purpose of Spring Profiles?}

Spring Profiles allow different configurations to be used for different
environments such as development, testing, and production.

For example:

\begin{verbatim}
application-dev.properties
application-test.properties
application-prod.properties
\end{verbatim}

A profile can be activated using:

\begin{verbatim}
spring.profiles.active=dev
\end{verbatim}

Profiles can also control which beans are created:

\begin{lstlisting}[language=Java]
@Configuration
@Profile("dev")
public class DevConfig {
}
\end{lstlisting}

This configuration is only loaded when the \texttt{dev} profile is active.

\subsubsection*{Interview Answer}

Spring Profiles allow environment-specific configuration. For example,
development, testing, and production environments can use different
database settings or beans without changing the application code.

\subsection{What is the Difference Between JPA and Hibernate?}

JPA is a specification that defines how Java objects should be persisted
and mapped to relational databases.

Hibernate is an implementation of the JPA specification.

\begin{verbatim}
JPA       -> specification / contract
Hibernate -> implementation
\end{verbatim}

For example, annotations such as \texttt{@Entity}, \texttt{@Id}, and
\texttt{@OneToMany} are defined by JPA, while Hibernate performs the actual
ORM work and communicates with the database.

Spring Data JPA provides another abstraction on top of JPA and simplifies
repository operations.

\begin{verbatim}
Spring Data JPA
      |
      v
     JPA
      |
      v
 Hibernate
      |
      v
 Database
\end{verbatim}

\subsubsection*{Interview Answer}

JPA is a specification for object-relational mapping in Java, while
Hibernate is one of its implementations. JPA defines the standard API and
rules, and Hibernate performs the actual persistence operations.

\subsection{How Do You Perform CRUD Operations Using JPA?}

In standard JPA, CRUD operations can be performed using
\texttt{EntityManager}:

\begin{verbatim}
Create -> persist()
Read   -> find()
Update -> merge()
Delete -> remove()
\end{verbatim}

For example:

\begin{lstlisting}[language=Java]
entityManager.persist(user);

User user =
    entityManager.find(User.class, 1L);

user.setName("Utku");
entityManager.merge(user);

entityManager.remove(user);
\end{lstlisting}

With Spring Data JPA, these operations are usually simplified through
\texttt{JpaRepository}:

\begin{lstlisting}[language=Java]
public interface UserRepository
        extends JpaRepository<User, Long> {
}
\end{lstlisting}

Common methods include:

\begin{verbatim}
save()       -> Create or Update
findById()   -> Read
findAll()    -> Read all
deleteById() -> Delete
\end{verbatim}

\subsubsection*{Interview Answer}

In standard JPA, CRUD operations are performed with EntityManager using
persist, find, merge, and remove. With Spring Data JPA, JpaRepository
provides methods such as save, findById, findAll, and deleteById.

\subsection{How Do You Implement Pagination in JPA?}

With Spring Data JPA, pagination is commonly implemented using
\texttt{Pageable} and \texttt{Page<T>}.

\begin{lstlisting}[language=Java]
Pageable pageable = PageRequest.of(0, 10);

Page<User> users =
        userRepository.findAll(pageable);
\end{lstlisting}

Here, \texttt{0} is the page number and \texttt{10} is the page size.

The returned \texttt{Page} contains both the current page data and
pagination metadata such as total pages and total elements.

Sorting can also be added:

\begin{lstlisting}[language=Java]
Pageable pageable = PageRequest.of(
        0,
        10,
        Sort.by("name").ascending()
);
\end{lstlisting}

In standard JPA, pagination can also be implemented using
\texttt{setFirstResult()} and \texttt{setMaxResults()}.

\subsubsection*{Interview Answer}

With Spring Data JPA, pagination is implemented using \texttt{Pageable},
usually created with \texttt{PageRequest}, and repository methods return a
\texttt{Page<T>} containing both the records and pagination metadata.

\subsection{What Are the Different Identifier Generation Strategies in JPA?}

JPA can generate primary key values using the
\texttt{@GeneratedValue} annotation.

Common strategies include:

\begin{verbatim}
IDENTITY
SEQUENCE
TABLE
AUTO
UUID
\end{verbatim}

For example:

\begin{lstlisting}[language=Java]
@Id
@GeneratedValue(strategy = GenerationType.IDENTITY)
private Long id;
\end{lstlisting}

\texttt{IDENTITY} uses the database's identity or auto-increment mechanism.

\texttt{SEQUENCE} uses a database sequence.

\texttt{TABLE} uses a separate table to generate identifiers.

\texttt{AUTO} allows the JPA provider to choose an appropriate strategy.

Modern JPA versions also support \texttt{UUID} generation.

\subsubsection*{Interview Answer}

JPA supports identifier generation strategies such as
\texttt{IDENTITY}, \texttt{SEQUENCE}, \texttt{TABLE},
\texttt{AUTO}, and \texttt{UUID}. The appropriate strategy depends on
the database and application requirements.

\subsection{What Fetch Strategies Are Available in JPA?}

JPA provides two main fetch strategies: \texttt{LAZY} and
\texttt{EAGER}.

Suppose a \texttt{User} has many orders:

\begin{lstlisting}[language=Java]
@OneToMany(fetch = FetchType.LAZY)
private List<Order> orders;
\end{lstlisting}

With \texttt{LAZY} loading, retrieving the user does not immediately load
the orders.

\begin{lstlisting}[language=Java]
User user = userRepository.findById(1L)
        .orElseThrow();

user.getName();   // Orders are not needed yet
user.getOrders(); // Orders may now be loaded
\end{lstlisting}

With \texttt{EAGER} loading:

\begin{lstlisting}[language=Java]
@OneToMany(fetch = FetchType.EAGER)
private List<Order> orders;
\end{lstlisting}

the related orders are loaded as part of loading the user, even if they are
not immediately used.

Default fetch types are:

\begin{verbatim}
@OneToMany  -> LAZY
@ManyToMany -> LAZY
@ManyToOne  -> EAGER
@OneToOne   -> EAGER
\end{verbatim}

\subsubsection*{Interview Answer}

LAZY loading delays loading related entities until they are needed, while
EAGER loading fetches them together with the main entity. LAZY loading can
avoid retrieving unnecessary data and is often preferred for large
relationships.

\subsection{What Caching Mechanisms Are Supported by JPA?}

JPA mainly provides two levels of caching.

\subsubsection*{First-Level Cache}

The first-level cache belongs to the persistence context and is associated
with an \texttt{EntityManager}.

It is automatic.

\begin{lstlisting}[language=Java]
User u1 = entityManager.find(User.class, 1L);
User u2 = entityManager.find(User.class, 1L);
\end{lstlisting}

Within the same persistence context, the second lookup can reuse the
already managed entity instead of querying the database again.

\subsubsection*{Second-Level Cache}

The second-level cache can be shared between multiple persistence contexts.

It is optional and depends on the JPA provider and its configuration.

\begin{verbatim}
First-Level Cache
-> EntityManager / Persistence Context scoped
-> Automatic

Second-Level Cache
-> Shared across persistence contexts
-> Optional and provider-dependent
\end{verbatim}

Query caching is not a standard JPA cache mechanism and may be provided by
implementations such as Hibernate.

\subsubsection*{Interview Answer}

JPA has a first-level cache associated with the persistence context, which
is automatic. It can also use an optional second-level cache shared across
multiple persistence contexts, depending on the JPA provider and
configuration.

\subsection{What is a Composite Key in SQL?}

A composite key is a key made of two or more columns that together uniquely
identify a row.

For example:

\begin{lstlisting}[language=SQL]
CREATE TABLE student_course (
    student_id BIGINT,
    course_id BIGINT,
    grade VARCHAR(2),

    PRIMARY KEY (student_id, course_id)
);
\end{lstlisting}

Neither \texttt{student\_id} nor \texttt{course\_id} must be unique by
itself, but their combination must be unique.

Composite keys are commonly used in junction tables that represent
many-to-many relationships.

\subsubsection*{Interview Answer}

A composite key is a key composed of multiple columns that together uniquely
identify a row. It is commonly used in junction tables, such as a
student-course relationship.

\subsection{When Do We Use an OUTER JOIN in SQL?}

An \texttt{OUTER JOIN} is used when unmatched rows should also be included
in the result.

For example:

\begin{lstlisting}[language=SQL]
SELECT u.name, o.id
FROM users u
LEFT JOIN orders o
    ON u.id = o.user_id;
\end{lstlisting}

A \texttt{LEFT JOIN} returns all rows from the left table and matching rows
from the right table.

If no matching row exists, the columns from the right table contain
\texttt{NULL}.

\begin{verbatim}
INNER JOIN -> only matching rows
LEFT JOIN  -> all rows from the left table
RIGHT JOIN -> all rows from the right table
FULL JOIN  -> all rows from both tables
\end{verbatim}

\subsubsection*{Interview Answer}

OUTER JOINs are used when we want to include unmatched rows in the result.
For example, a LEFT JOIN can return all users even if some users have no
orders.

\subsection{What is a Database View and When Do We Use It?}

A database \texttt{VIEW} is a virtual table based on a SQL query.

Views are commonly used to simplify complex queries, reuse query logic,
restrict access to specific columns or rows, and provide abstraction over
underlying tables.

A normal view generally stores the query definition rather than storing the
query result itself.

\subsubsection*{Interview Answer}

A database view is a virtual table based on a SQL query. It is useful for
simplifying complex queries, reusing query logic, restricting access to
specific data, and hiding underlying table complexity.

\subsection{What is the Difference Between a Process and a Thread?}

A process is an independent running program with its own memory space.

A thread is a unit of execution inside a process.

Multiple threads in the same process share resources such as heap memory,
while each thread has its own stack.

\begin{verbatim}
Process
|
+-- Shared Heap
|
+-- Thread 1 -> own stack
+-- Thread 2 -> own stack
+-- Thread 3 -> own stack
\end{verbatim}

Processes are more isolated and generally more expensive to create, while
threads are lighter and can communicate through shared memory.

However, shared memory can introduce concurrency problems such as race
conditions.

\subsubsection*{Interview Answer}

A process is an independent running program with its own memory space,
while a thread is a unit of execution inside a process. Threads in the same
process share resources such as heap memory, but each thread has its own
stack.

\subsection{What is Multithreading?}

Multithreading means executing multiple threads within the same process.

Threads share process resources such as heap memory, while each thread has
its own stack.



Calling \texttt{start()} creates a new thread of execution, while calling
\texttt{run()} directly is only a normal method call.

Multithreading can improve responsiveness and concurrency, but shared state
must be handled carefully to avoid problems such as race conditions.

\subsubsection*{Interview Answer}

Multithreading is the execution of multiple threads within the same process.
Threads share process resources but have separate stacks. It improves
concurrency, but shared mutable state can introduce thread-safety problems.

\subsection{What is Multiprocessing?}

Multiprocessing means running multiple independent processes at the same
time.

Each process has its own memory space and resources.

\begin{verbatim}
Process 1
-> Separate memory

Process 2
-> Separate memory
\end{verbatim}

This is different from multithreading, where multiple threads run inside
the same process and share memory.

Multiprocessing provides better isolation and can make effective use of
multiple CPU cores.

However, communication between processes is generally more expensive and
usually requires mechanisms such as pipes, sockets, shared memory, or
message queues.

\subsubsection*{Interview Answer}

Multiprocessing means running multiple independent processes, each with its
own memory space. It provides strong isolation and can use multiple CPU
cores effectively, but inter-process communication is more expensive than
communication between threads.

\subsection{What is the Difference Between Multithreading and Multiprocessing?}

Multithreading uses multiple threads inside the same process, while
multiprocessing uses multiple independent processes.

\begin{verbatim}
Multithreading
-> One process
-> Multiple threads
-> Shared memory
-> Lower overhead

Multiprocessing
-> Multiple processes
-> Separate memory
-> Better isolation
-> Higher overhead
\end{verbatim}

Threads can communicate easily through shared memory, but this may create
race conditions and synchronization problems.

Processes are more isolated, but communication between them usually
requires inter-process communication mechanisms.

\subsubsection*{Interview Answer}

Multithreading uses multiple threads inside the same process, so they share
memory and are lightweight. Multiprocessing uses separate processes with
independent memory spaces, providing better isolation but higher overhead.

\subsection{What is a Race Condition?}

A race condition occurs when multiple threads access and modify shared
mutable data concurrently, causing the result to depend on the timing of
the threads.

For example:

\begin{lstlisting}[language=Java]
class Counter {

    private int count = 0;

    public void increment() {
        count++;
    }
}
\end{lstlisting}

If multiple threads call \texttt{increment()} at the same time, some
updates may be lost because \texttt{count++} is not an atomic operation.

Race conditions can be prevented using synchronization, locks, or atomic
classes such as \texttt{AtomicInteger}.

\subsubsection*{Interview Answer}

A race condition occurs when multiple threads modify shared data
concurrently and the result depends on the order or timing of execution.
Proper synchronization or atomic operations can be used to prevent it.

\subsection{What is Synchronization and Why is it Needed?}

Synchronization controls access to shared mutable data when multiple
threads are running concurrently.

For example:

\begin{lstlisting}[language=Java]
class Counter {

    private int count = 0;

    public synchronized void increment() {
        count++;
    }
}
\end{lstlisting}

The \texttt{synchronized} keyword ensures that only one thread at a time
can execute the protected section for the same object.

Synchronization helps prevent race conditions and inconsistent data.

However, excessive synchronization can reduce performance because threads
may need to wait for locks.

\subsubsection*{Interview Answer}

Synchronization controls concurrent access to shared mutable data. It
protects critical sections so that only one thread at a time can execute
them, helping prevent race conditions.

\subsection{What is a Deadlock?}

A deadlock occurs when two or more threads wait indefinitely for resources
held by each other, so none of them can continue execution.

For example:

\begin{verbatim}
Thread 1:
holds Lock A
waits for Lock B

Thread 2:
holds Lock B
waits for Lock A
\end{verbatim}

This creates a circular dependency.

One common way to reduce deadlock risk is to acquire locks in a consistent
order.

\subsubsection*{Interview Answer}

A deadlock occurs when threads wait indefinitely for locks held by each
other, preventing all of them from making progress.

\subsection{What Does Thread-Safe Mean?}

Thread-safe code behaves correctly when it is accessed by multiple threads
concurrently.

For example, this is not thread-safe:

\begin{lstlisting}[language=Java]
class Counter {
    private int count = 0;

    public void increment() {
        count++;
    }
}
\end{lstlisting}

It can be made thread-safe using synchronization:

\begin{lstlisting}[language=Java]
public synchronized void increment() {
    count++;
}
\end{lstlisting}

Other approaches include atomic classes, locks, and immutable objects.

\subsubsection*{Interview Answer}

Thread-safe code works correctly when multiple threads access it
concurrently. Shared mutable state must be protected from race conditions
and inconsistent updates.

\subsection{What Are the Advantages of Stored Procedures?}

A stored procedure is a reusable set of SQL statements stored and executed
inside the database.

Main advantages include:

\begin{itemize}
    \item Reusing database logic
    \item Reducing network communication
    \item Centralizing complex database operations
    \item Providing controlled access to data
    \item Potential performance benefits depending on the DBMS
\end{itemize}

\subsubsection*{Interview Answer}

A stored procedure is a reusable set of SQL statements stored in the
database. It can reduce network communication, centralize database logic,
improve access control, and provide performance benefits in some cases.
\end{document}